% Architecture en couches — bandes contiguës, teinte croissante vers le bas
\begin{tikzpicture}[
  band/.style={base, rectangle, minimum width=\textwidth, minimum height=1.25cm,
    inner xsep=0pt, align=left, text width=\textwidth-3.6cm, anchor=west}
]
  \def\w{\textwidth}
  \foreach \n/\t/\d/\f in {
    5/{Services utilisateurs}/{IaaS, cyberlabs, services communs}/4,
    4/{Orchestration}/{OpenStack / Proxmox, quotas, gouvernance}/8,
    3/{Hyperviseurs}/{Proxmox (transitoire), OpenStack (cible)}/12,
    2/{Réseau}/{SRX300, switches Cisco, VLAN, VPN}/16,
    1/{Physique}/{serveurs, stockage, alimentation, climatisation}/20}
  {
    \pgfmathsetmacro{\y}{-(5-\n)*1.25}
    \node[base, rectangle, fill=accent!\f, minimum width=\textwidth, minimum height=1.25cm,
      anchor=north west, inner sep=0pt] (b\n) at (0,\y) {};
    \node[font=\sffamily\scriptsize\bfseries, text=accent, anchor=west] at ([xshift=10pt]b\n.west)
      {\MakeUppercase{Couche \n}};
    \node[font=\sffamily\footnotesize, anchor=west, text=ink] at ([xshift=2.5cm]b\n.west)
      {\textbf{\t}\ \ \textcolor{muted}{—}\ \ \d};
  }
\end{tikzpicture}
